(1) 自旋二分之一粒子的哈密顿量为 \(H(t)=-\mu B\,\boldsymbol n(t)\cdot\boldsymbol\sigma\)，\(\mu B>0\) 恒定，磁场单位矢量 \(\boldsymbol n\) 取固定极角 \(\vartheta\)、方位角 \(\phi(t)\) 从0缓慢增加至 \(2\pi\)。求瞬时能级及能隙。
(2) 给出沿场低能本征态的一种规范化自旋量子态，并求 Berry 连接 \(i\langle +|\partial_\phi|+\rangle\)。
(3) 求一周的几何相位和动力学相位，并将几何相位写成扫过的固角。
(4) 讨论 \(\vartheta\to0\)、\(\vartheta\to\pi/2\) 两极限以及绝热条件；说明相位为何可在干涉实验中观测。
